% Finished replacement passages spanning the whole argument.
% This is a review proposal, not an edit to any frozen full draft.
% Quote macros come unchanged from the shared manifest, Opus's supplement,
% and Fable 5.2's supplement. Compile through replacement_preview.tex.

\title{Bargaining, development and creditor inheritance in MonopolyBench}
\maketitle
\begin{abstract}
Economic agents trade with competitors whose construction budgets and survival affect their own positions. We study these interactions in twenty complete four-player MonopolyBench games involving four configured LLM systems. Across 918 acceptance-legal offers, responses differ with the ownership consequences and negotiated terms. Counteroffers often preserve a rival's proposed completion while changing compensation: 92 of 111 counters to author-only completing offers retain that same completion. Explicit explanations and contextual exchanges discuss rival development and prospective creditor inheritance, although neither cash direction nor rising counteroffers identify a completion premium. Development also depends on access and execution. Of 63 trade-completed groups, 43 receive construction in the first management window; other windows contain unavailable options, defaults, conflicting channels or valid ends. Retained local financing checks distinguish three payable elimination debts, fifty-six beyond a generous own-asset bound and one hotel-stock lock. Estates complete 52 groups while play continues, with another 69 completions arriving only at game end. Eight contextual incidents connect these populations to bargaining, correction and execution. The findings describe economic interaction under the cohort's unequal response policies and inconsistent history summaries. They establish neither optimal strategy nor causal persuasion, durable learning, a general capability ranking or an effect of inheritance on winning.
\end{abstract}

\section{Introduction}

A property exchange changes more than the recipient's holdings. It can finance the seller's construction, remove a blocker against another player, or supply liquidity to someone whose bankruptcy would strengthen a creditor. The value of a trade therefore depends on what both parties can do afterward and on the assets that might pass to others if either fails.

We study this problem in twenty complete MonopolyBench games. The four configured agents bargain over deeds and cash, invest in houses and hotels, liquidate assets and acquire estates. Every offer has recorded terms, every decision has an engine-issued menu, and each applied action has resulting events. Those records allow us to follow economic interactions while distinguishing a stated reason, a requested action and what the system actually executed.

Our central finding is a connected description of economic interdependence. Some agents explicitly price or refuse a completion because of the rival's resulting building capacity. Some discuss preserving another player's solvency to deny its estate to a third party. These considerations appear alongside differentiated response patterns in the complete negotiation record. The acquired groups then follow unequal paths into construction and receipts, and bankruptcy redistributes deeds and sometimes buildings to creditors. We report that sequence without assigning optimality to an observed choice or causal credit to a message.

The evaluated objects are complete configurations. Provider refusals change one participant's response policy and memory, historical summaries sometimes misdescribe trades, and some received tool calls conflict with accompanying text. These are conditions of the observed behavior. Reconstruction and retained legal witnesses help specify what happened under them; they do not produce a clean comparison of the underlying models.

\section{Environment and observation conditions}

MonopolyBench uses the forty-space United States board with four players, \$1,500 starting cash each, 32 bank houses and 12 hotels. The deterministic engine alone changes state, emits legal actions and records events. Declining a purchase starts an ascending auction. Bilateral trade threads permit repeated counters before normal play resumes. Deeds in developed groups cannot be traded, and recipients of mortgaged deeds pay transfer interest.

Several implemented rules shape the economic analysis. Construction requires complete, unmortgaged ownership, cash and stock, and each build-plan entry must preserve even building. A player in liquidation can sell buildings, mortgage eligible deeds or declare bankruptcy, but cannot start a rescue trade. Removing buildings can unlock mortgages. Hotel downgrades require bank houses. A declaration remains available when another own-asset sequence could pay. Player-creditors inherit buildings intact, while bank-returned deeds are not immediately auctioned. The implementation's floating-point redemption charge can exceed the intended charge by one dollar. Our historical claims concern this preserved implementation.

Five seeds are each played in four cyclic seat rotations, yielding twenty games, eighty player-games, 6,989 decisions, 8,679 attempts and 50,918 events. All games end by bankruptcy without a turn cap. Each configuration occupies every absolute seat five times and faces the same three opponents. These are dependent trajectories, not independent decisions sampled for capability inference. Shared seeds need not preserve later landings after actions change random-draw consumption.

The participants use recorded OpenRouter identifiers \path{openai/gpt-6.1-sol}, \path{google/gemini-4-argon}, \path{anthropic/claude-opus-5.5} and \path{x-ai/grok-4.7}. All have recorded medium reasoning effort, which does not equalize computation. Grok uses required tool choice. Invalid output ordinarily receives a corrective request within a two-attempt budget. Opus has 1,604 first-attempt refusals followed by action-only recovery; one further refusal follows an invalid first attempt at G2 d129. Recovery removes the private-note field and consumes the remaining attempt. In 1,394 Opus decisions, its private-memory window contains no model-authored note. The refusal trigger and its change across the source boundary remain unresolved.

Agents receive public cash, ownership and development, visible model names, bounded public and own-private memory, and structured negotiation history. A formatter confuses current offer authors with thread initiators. Known manifestations occur in 2,054 initial prompts, despite correct cash and holdings. Current-offer acceptance accounting nevertheless agrees with the delivered fields in all 920 response opportunities. Public logging is also distinct from delivery: 149 of 6,882 nonempty message occurrences never enter the timeline. Focal reception claims rely on actual recipient prompts.

Grok's selected records contain 139 content/tool action-name conflicts already present in saved nonstreamed received responses. That excludes local streaming assembly for these records; generation, provider processing and gateway behavior remain unresolved. Sol and Gemini's empty selected content fields prevent the same comparison. The engine applies the received valid call.

Standard prompts permit public deception and request private notes reflecting true reasoning. Action-only recovery removes the private-note field and that honesty clause. A false public statement is not by itself an instruction violation. A false private note can express a mistaken belief and does not alone establish dishonesty. None of the recorded text channels is direct access to internal computation.

\section{Measures and evidence}

A response opportunity is one received offer. A counter changes the author of the current terms, so author and original initiator must remain distinct. Color completion is an ownership fact, including mortgaged deeds and excluding railroads and utilities. It does not imply buildability or utility. Cash direction describes the negotiated transfer, with mandatory bank charges accounted for separately.

The first management window contains the owner's actual management decisions in its first management turn after a completion. We separately record construction available on entry, construction appearing later on recorded menus, and applied building. Absence from those menus does not prove that every unvisited financing sequence fails. A spending unit is a game/player/color triple with positive own house or hotel expenditure. Developed-rent receipts differ from all rent received on that color, and neither is incremental investment return. Acquisition, redemption, salvage and estate flows remain separate.

Retained legal checks fix a historical prefix, actor, permitted own actions and endpoint. Construction checks preserve requested purchases while splitting count-valued entries and ordering them legally. Current-debt checks permit own sales and mortgages. For cash $C$, full nominal salvage $S$ and mortgage capacity $M$ including collateral unlocked by sales, $U=C+S+M$ is a generous bound. If $U$ is below the debt, own-asset payment is impossible. Otherwise a positive witness or a suitably restricted negative search is required. These endpoints establish local feasibility, not sound investment, future survival or a different winner.

Historical reconstruction matches all initial decision states and saved snapshots, and retained action streams replay under the preserved engine. Those checks support implementation consistency. They do not independently establish intended-rule correctness or semantic validity. The current analysis re-tabulates existing ledgers and checks original quote fields and selected delivered prompts; it does not newly execute the engine certificates.

Semantic evidence is explicitly LLM-assisted. Two Astra research instances read seventy-four rejections completing the recipient's group, with partial prior exposure for the second reader; one reads the temporal sample. Other extraction and review uses model families overlapping evaluated participants. Agreement among them is not human annotation or an independent validity standard. The eight incidents below are purposively selected for distinct economic mechanisms and evidential contrasts, not prevalence.

\section{Bargaining over development opportunities}

There are 526 negotiation threads and 920 responses, comprising 113 acceptances, 394 counters and 413 rejections. Acceptance is legal in 918 cases. None resolves through an engine default, although 168 responses use a retry. Author-only completing offers receive 19 acceptances, 111 counters and 159 rejections in 289 opportunities. Recipient-only offers receive 24, 78 and 38 in 140. The median negotiated transfer to the recipient is \$200 in the former cell and minus \$350 in the latter. The completing side pays in 373 of 429 one-sided offers. These are selected, heterogeneous bundles, not an estimate of the price or causal effect of completion.

Configuration and thread position matter. Sol, Gemini and Grok accept author-only completion less often than recipient-only completion. Opus reverses that ordering, accepting 7/52 and 5/42 respectively and usually countering rather than rejecting. Author-only offers are also more often opening bids. Within first and later offers, the rejection contrast retains its direction but shrinks. Repeated offers within games and threads remain dependent.

Counter content supplies a more specific observation. Ninety-two of 111 counters to author-only offers preserve the same original author's newly complete group. Eighty-eight leave only that party completing; 85 improve the counterer's net cash, and 79 preserve the entire noncash bundle while doing so. These are negotiations over compensation with the completing deed still available. Cash demands also rise in other offer categories, so this does not identify a completion premium. Source-attributed reasons supply the separate evidence that some agents expressly value rival consequences. Both LLM readings identify such a stated consideration in 61 of the seventy-four recipient-completion rejections, without establishing the quality or cause of those refusals.

Seller cash differs among settlements. In 48 cash settlements completing the payer's group, median seller cash is \$194, compared with \$571 in 57 cash settlements completing neither party's group. Twenty-one of the 48 sellers hold less than \$100. That is an exploratory association, not a majority of completing sales or evidence that scarcity caused release.

\TranscriptTitle{sol-g2}{A completion refused and later traded in reverse}

At G2 turn 45, Opus holds two yellows and offers Sol \$325 for Marvin Gardens. Sol answers:
\TranscriptQuote{g2-d124-price}
\TranscriptQuote{g2-d124-note}
The residual cash figures are correct, but the house calculation is not. Nine yellow houses cost \$1,350; \$520 alone buys three. This is explicit attention to rival development with a mistaken estimate of its immediate cash-based scale.

Opus later calls \$450 final while recording a private \$475 ceiling. Sol refuses. At turn 69 the exchange reverses: Sol acquires Atlantic and Ventnor for two railroads and \$100. It spends \$300 on yellow houses and recovers \$150 through sales. It receives no developed yellow rent, but does receive \$44 of undeveloped Atlantic rent at turn 77. Its later bankruptcy transfers the deeds to Gemini, the eventual winner. The sequence connects negotiation, investment, liquidation and creditor acquisition; it does not establish which price or trade was best.
\TranscriptSource{G2 d124--128 and d197--198; received public messages checked in subsequent counterparty requests. The \$44 receipt is original event 1566. The developed-rent measure is separate.}

\TranscriptTitle{sol-g15}{Both sides discuss asymmetric construction cash}

In G15's four-player state, Sol offers two greens for Grok's Ventnor and \$100, completing yellow for Sol and green for Grok. Their post-acceptance cash would be \$830 and \$233. Grok explains:
\TranscriptQuote{fx-g15-refusal-note}
The \$600 figure is the cost of one house on each green, not the minimum legal construction. A single green house costs \$200. The note therefore distinguishes capacity to fund a full initial layer, while leaving collateral and later choices unestablished.

Sol drops its cash request and adds Kentucky, then offers both greens, Kentucky and \$250. Grok still refuses:
\TranscriptQuote{fx-g15-final-refusal}
No trade occurs. Eighteen turns later Grok is eliminated first and its blockers pass to Gemini, which wins. The case supplies direct rival-development reasoning, with both parties addressing the same cash asymmetry. It does not prove optimal refusal or the maximum construction each could finance.
\TranscriptSource{G15 d168--175, standard first attempts. New public-message spans are present in the actual subsequent counterparty requests.}

\TranscriptTitle{sol-g7}{A settled bluff about a final limit}

Grok wants North Carolina and publicly offers:
\TranscriptQuote{g7_d42_public}
Its applied private note retains a higher contingency:
\TranscriptQuote{g7_d42_private}
Sol accepts \$600 and states a preference for reinvesting over seeking another \$25. The public/private discrepancy is a clear authorized bargaining bluff, and contrasts with G2's unaccepted final offer. The private contingency is not tested, and the record does not identify whether ``final'' changed the price.
\TranscriptSource{G7 d41--43, standard first attempts; the \$600 message is delivered to Sol.}

\TranscriptTitle{sol-g19}{Correction accompanies acceptance}

Grok's received proposal offers Virginia plus \$150 for States, falsely presenting States as an orange deed. Its content selects END and its attached note says to abort. Sol replies:
\TranscriptQuote{g19_d136_correction_accept}
The swap and payment execute. This is evidence of explicit rejection of the factual premise with acceptance of the terms. The proposal's public message agrees with the received call; its content and attached note oppose it. These conflicting channels do not establish one coherent plan to deceive.
\TranscriptSource{G19 d135--136, sole valid standard responses.}

\section{From ownership to executed development}

All 63 trade-completed groups reach a first management window. Forty-three receive construction. The twenty residual windows contain nine with no construction option on any recorded menu, four engine defaults, four received END conflicts and three valid-end windows. Among the 42 buildable on entry, 34 receive construction; the eight residuals contain three defaults, three conflicts and two stated holds. The populations describe access at different times and must not be combined into a willingness-to-invest rate.

Seven of the nine no-option groups receive a mortgaged completing deed; two settlements leave less cash than one house costs. These describe immediate constraints. They do not exhaust unvisited financing routes, and stated intentions to wait can coexist with an absent option. Likewise, a valid END without a detected conflict is not automatically a verified reserve preference. Genuine reserve explanations occur in G3 d176 and G1 d426, without an established optimality claim.

The retained audit finds all 61 defaulted Opus purchase multisets reachable after atomic expansion and legal ordering, while their submitted encodings remain invalid. Sixty chains have refusal before invalid action-only construction; G2 d129 has the reverse order. For Opus's first-window nonbuilds specifically, five lack a recorded construction option and three end in defaults. Feasibility and failed execution are demonstrated at the stated endpoints. Neither the amount of lost rent nor the investment quality of the targets is measured.

\TranscriptTitle{sol-g20}{An appeal to withhold rival construction cash}

After Grok buys Illinois for \$500 and completes reds, Gemini tells the table:
\TranscriptQuote{g20-coordination-appeal}
Grok's next note recognizes a boycott. It later offers Vermont to Gemini for \$150, and Gemini refuses:
\TranscriptQuote{g20-concrete-refusal}
Grok mortgages Vermont instead. The target recognizes the appeal and its author follows it in a concrete response. Sol had already expressed its own preference not to finance Grok, and Opus subsequently proposes cash purchases. This is attempted financial denial with unilateral follow-through, not evidence of a jointly adopted cartel or a changed policy caused by the appeal.
\TranscriptSource{G20 d95--102. The appeal reaches Grok's d97 prompt; the note is from selected attempt 1.}

\section{Current-debt finance and creditor redistribution}

Fifty-six of sixty elimination debts exceed the generous own-asset bound. Three have retained legal payment witnesses. The remaining case, G7 d339, has \$1,310 nominal resources but only \$440 reachable against \$750 owed. With no bank houses, its pink hotel cannot be downgraded and even-selling blocks a lower sibling's house. The mortgage-only continuation triggers automatic bankruptcy. This is a local financing impossibility under the preserved rules, not a judgment about prior investment.

The payable cases have different explanations. Sol's G12 d190 and Gemini's G17 d234 accounts omit collateral unlocked by selling buildings. Grok's G12 d249 instead contains conflicting channels. A mortgage-only subtotal is not the full displayed menu, which includes sales. Actual staged finance also succeeds, as G20 d327--336 uses \$49 cash, \$700 salvage and \$665 mortgages to pay \$1,400 and retain \$14.

\TranscriptTitle{sol-g12}{A payable debt ends in contradictory execution}

At G12 d249, Grok owes \$250 with \$43 cash. Its content requests a two-house Virginia sale. Its note concludes:
\TranscriptQuote{g12_d249_private_final}
The sole received tool declares bankruptcy, with the public message:
\TranscriptQuote{g12_d249_public_bankruptcy}
The engine eliminates Grok without a sale. A retained five-sale witness pays and leaves \$43. The requested two-house sale is legal but insufficient, reaching \$143. Content and note oppose declaration; the public message supports it. The evidence establishes a current-debt execution failure with unresolved origin, rather than deliberate surrender or a unified mortgage-omission error.
\TranscriptSource{G12 d249, standard attempt 0. The payment witness is inherited from the historical audit.}

During continuing play, estates complete 52 color groups, compared with 63 trade completions. Another 69 estate completions occur at terminal bankruptcy and supply no later management opportunity. Eventual winners receive 202 of 227 nonterminal inherited deeds and are creditors in 35 of forty nonterminal eliminations. Those outcome-conditioned counts describe concentration, not an effect of inheritance on winning. The rule makes solvency and creditor identity economically relevant, as the next exchanges explicitly show.

\TranscriptTitle{sol-g13}{An insurance pitch omits collateral}

Grok opens G13 d250 with \$35, offering mortgaged Pennsylvania for \$200 and warning that Short Line rent would transfer yellows to Sol. Opus's \$130 counter and solvency rationale follow that framing. Grok replies:
\TranscriptQuote{g13-seller-counter}
Its returned provider explanation states:
\TranscriptQuote{g13-omission-explanation}
Its note separately recognizes Ventnor collateral and sets a \$165 floor. Opus accepts \$200 and pays \$16 bank interest. The retained conditional \$130-acceptance route returns management to Grok, where a \$130 Ventnor mortgage reaches \$295. That establishes a cheaper available financing route, not that Grok would accept it.

The \$70 difference is observed; its cause is not. Opus echoes estate-denial stakes without endorsing the alleged minimum. Grok may also be expressing refusal to transact at \$130. The longer provider explanation contradicts itself, and its opening 5/36 probability is wrong: Short Line is 4/36 from Atlantic. Forty-three turns after settlement, Grok fails on Short Line and yellows pass to Sol. That chronology does not establish forty-three turns of purchased survival.
\TranscriptSource{G13 d250--253 and d311. Opus uses action-only recovery; focal delivered histories contain self-partner wording.}

\TranscriptTitle{sol-g18}{A standing price is withdrawn after deterioration}

Opus repeatedly offers \$250 for Sol's mortgaged Pacific and, after collecting \$750 rent, says:
\TranscriptQuote{g18-assurance}
A further \$750 rent forces Sol to liquidate, leaving \$49. Sol invokes the advertised offer and receives:
\TranscriptQuote{g18-invocation}
\TranscriptQuote{g18-withdrawal}
The returned explanation cites Sol's near-insolvency and Opus's redemption and building costs. No sale occurs. The assurance is delivered at invocation and withdrawal, so its absence from the observed request does not explain the reversal. The episode records changed terms and an explicit financing rationale. It establishes neither original insincerity nor costly reliance, and the engine supplies no binding standing-offer object or expiry.
\TranscriptSource{G18 d421--442; Opus's explanation is provider-returned text under action-only recovery.}

\section{Plans, revision and the limits of the broader questions}

The fixed 240-decision text sample contains particular targets, forecasts and revisions rather than a planning score. Six of eight named landing-purchase policies never meet a later opportunity. One is fulfilled; G16's Virginia policy meets an affordable purchase menu but resolves to an auction despite a note favoring purchase. Conversely, Sol's G9 target of three houses on each dark blue is realized over 31 turns. The record therefore contains both sustained targets and conflicting execution, without identifying a memory benefit.

Construction errors also recur under unequal exposure. Of 114 uneven-error decisions, 109 have a later submission and 63 repeat the error. Among fifty next submissions after a decision containing a corrective request with the exact error, nine repeat it. Opus's refusal path often consumes that request opportunity. These are observed sequences under changing states, not a feedback effect or durable learning estimate.

The evidence receives different levels for the motivating questions. Q1 strategic reasoning has explicit considerations and local executable alternatives, not optimality. Q2 opponent modeling has selected public-capacity deductions, not calibrated hidden beliefs. Q3 negotiation and conduct has a complete response census and contextual tactics, not identified efficacy, fairness, agreement bias or prevalence. Q4 horizon and memory has targets and actual opportunities, not a causal memory benefit. Q5 adaptation has revisions and repair, not durable learning or unfamiliar-opponent adaptation. Q6 economics has actual flows and current-debt endpoints, not ex ante values or risk-adjusted returns. Q7 collective dynamics has creditor transfers and attempted coordination, not coalition stability or collusion. Q8 validity has reconstruction under stated conditions, not a general ranking or transfer result.

For the broader questions, A receives concrete cases where talk and action correspond or diverge, not a general competence verdict. B, the relation between general and multi-agent capability, is untested. C receives specific error and inheritance pathways, not a decomposition of wins into skill, luck and opponent weakness.

\section{Relation to earlier studies}

Economic-agent studies already connect process and outcomes. AgentMonopoly measures actions, language and persona interventions in Monopoly; its v2 prompt permits bluffing while prohibiting messages contradicting the chosen action~\citep{wu2026agentmonopoly}. Arena v3 analyzes state-grounded inter-agent communication and repeats its complete pipeline with replacement judge families~\citep{li2026arena}. Our selected cases do not provide its population communication rate. FinBoardBench is precedent for recording response intervention in financial board games, while DECE uses matched execution-framework comparisons that this cohort lacks~\citep{hu2026finboardbench,zhang2026dece}. NegotiationArena's controlled tactic changes provide a different basis for efficacy claims than our observed settlements~\citep{bianchi2024negotiationarena}. The present contribution is the connected finite-cohort account of bargaining, executable development, current-debt finance and creditor succession. It claims no novelty for Monopoly agents, failure logging, replay or state-grounded communication analysis.

\section{Limits and research provenance}

The five seed blocks, fixed roster, visible names, response conditions and implemented rules define the domain. Retrospective semantic selection, uneven channels, LLM-family overlap and partial exposure limit interpretation. The 126 fixed historical cases support declared endpoints; they are not 126 independent mechanisms or a general semantic-verifier evaluation. The negative checker does not certify arbitrary bundled-sale impossibility. No new games, model evaluations or engine alternatives were executed for this analysis.

Human-directed research assistance used Astra, Sol, Opus and Fable instances. Earlier work includes Fable 5.1; this rewrite uses complete independent papers by Astra, Sol, Opus and Fable 5.2 before cross-review. These are research roles, not additional gameplay participants or independent human annotation. The evidence archive preserves source identities, exact quote fields, disagreement, corrections and performed-versus-inherited checks. The historical checkpoint and the later rewrite must remain separately identified.

\section{Conclusion}

The cohort records agents negotiating over rival development and prospective estates, followed by different paths from ownership into construction, finance and creditor transfer. Its evidence includes explicit strategic considerations, competent factual correction, reachable allocations that fail execution, and both payable and inaccessible debts. These observations explain the economic sequence in the recorded systems. Their scope remains that sequence, rather than a causal account of persuasion, optimal investment or victory.
